\section{Related Work}
\label{sec:related}
\subsection{Alignment and Panorama Composition}
% Role: prior methods; establish the global-homography scope.
Brown and Lowe describe automatic panoramic stitching using invariant features
and multi-band blending~\cite{brown2007stitching}. Their work supplies the classical
pipeline context; this study concerns computational budgets within a GPU
implementation. Their feature matching is followed by RANSAC homography
estimation and inlier-based image verification. This establishes the distinction
between descriptor similarity and geometric consistency; it does not imply that
our implementation uses their SIFT features or nearest-neighbour search.
More flexible image models address a different limitation:
UDIS++ combines global homography with local thin-plate-spline warping and
learned seam-driven composition~\cite{nie2023udispp}. Our global homography
renderer does not reproduce those local deformations. Increased rendering cost
cannot be assumed to repair parallax or an incorrect alignment.

\subsection{Blending and GPU Cost--Quality Trade-offs}
% Role: foundations and evaluation precedent.
Multi-resolution splines blend band-pass components across spatial scales before
reconstruction~\cite{burt1983spline}. Our multi-band mode uses the
Gaussian/Laplacian-pyramid principle while fixing its own ownership, extension,
border, and sampling conventions. These choices require an explicit reference
rather than an assumption of bitwise equivalence to a library blender. Zhu et al.
compare panoramic-video blending algorithms using time, memory, and visual-quality
measurements, including GPU implementations where
appropriate~\cite{zhu2018blending}. Their evaluation motivates reporting these
costs separately. Their video measurements are not performance predictions for
our two-image workload or GPU platform.

\subsection{Configuration Tuning and Quality-Aware Execution}
% Role: acknowledge systems precedents without a first-of-kind claim.
Program autotuning treats implementation and parameter choices as an explicit
search space. OpenTuner supports domain-specific representations and multiple
search techniques for that purpose~\cite{ansel2014opentuner}. SAGE combines
compiler-generated approximate CUDA kernels with runtime selection under a user
quality target~\cite{samadi2013sage}. These systems establish configuration
search and quality-aware GPU execution as prior ideas. Our evaluation must
isolate scene-dependent budget selection, intermediate geometry checks, and
valid result reuse in the stitching pipeline. A comparison against fixed
budgets alone does not explain which runtime mechanism produces an improvement.
